\section{The Centroidal Voronoi Districts Algorithm}
\label{sec:algorithm}

\subsection{Notation and Subgraph Setup}

Let $G_v = (V_v, E_v, w_v)$ denote the subgraph at bisection tree node $v$,
where $V_v$ is the set of census tracts assigned to the subtree rooted at $v$,
$E_v \subseteq E$ is the induced edge set, and $w_v$ inherits weights from the
parent graph.
Let $m = |V_v|$ be the number of tracts.
Let $p_i$ denote the population of tract $i \in V_v$, and let
$P_v = \sum_{i \in V_v} p_i$ be the total subgraph population.

For bisection ($k=2$), \cvd{} assigns each tract to one of two districts
$\{D_0, D_1\}$ and returns a partition $(S, V_v \setminus S)$ satisfying
the population-balance constraint:
\[
  \left|\sum_{i \in S} p_i - \frac{k_L}{k_v} P_v \right|
  \;\le\;
  \delta \cdot \frac{P_v}{k_v},
\]
where $\delta$ is the balance tolerance (default $0.005$, the $\pm 0.5\%$
federal tolerance) and $k_L / k_v$ is the left-subtree population target.

Let $d_{\mathrm{BFS}}(i, j)$ denote the BFS hop-count distance between
tracts $i$ and $j$ in $G_v$: the length of the shortest path in $E_v$.
For non-adjacent tracts with no path, $d_{\mathrm{BFS}}(i,j) = \infty$;
this does not arise for connected subgraphs.

\subsection{Seed Derivation}

\cvd{} requires a per-node random seed to ensure reproducibility.
The seed must be deterministic given the global base seed and the node's
position in the bisection tree.

\begin{definition}[CVD Node Seed]
\label{def:cvd-seed}
Let $\mathtt{base\_seed}$ be the unsigned 64-bit integer global seed, and let
$\mathtt{node\_path}$ be the string representation of the path from the root
to node $v$ (e.g., \texttt{"0/1/0"} for the left-right-left child of the root).
The \emph{CVD node seed} is
\[
  s_v = \mathrm{first8}\bigl(
    \mathrm{SHA\text{-}256}(
      \texttt{"CVD\_INIT\_"} \;\|\; \mathtt{node\_path}
      \;\|\; \texttt{"\_"} \;\|\; \mathrm{le8}(\mathtt{base\_seed})
    )
  \bigr),
\]
where $\mathrm{le8}(x)$ encodes $x$ as 8 little-endian bytes and
$\mathrm{first8}(h)$ reads the first 8 bytes of hash $h$ as a
little-endian unsigned 64-bit integer~\citep{sha256nist}.
\end{definition}

The domain-separation prefix \texttt{"CVD\_INIT\_"} ensures that the CVD
seed is independent from all other SHA-256 applications in the platform
(e.g., \texttt{"SA\_NODE\_"} for B.19's simulated annealing).
A platform invariant test asserts this prefix constant equals
\texttt{"CVD\_INIT\_"} exactly, and that $s_v(\mathtt{base\_seed}=0)$
equals a hard-coded expected value.

\subsection{K-Farthest Seed Initialisation}

The two initial seeds are placed by the k-farthest BFS procedure:

\begin{definition}[K-Farthest Initialisation]
\label{def:kfarthest}
Let the tracts in $V_v$ be sorted by global index, giving an ordered list
$\mathrm{sorted} = (t_0, t_1, \ldots, t_{m-1})$.
\begin{enumerate}
\item \textbf{First seed}: $\mathrm{seed}_0 = t_{s_v \bmod m}$, where $s_v$ is
  the CVD node seed.
\item \textbf{BFS from $\mathrm{seed}_0$}: compute $d_0(t) = d_{\mathrm{BFS}}(\mathrm{seed}_0, t)$
  for all $t \in V_v$.
\item \textbf{Second seed}: $\mathrm{seed}_1 = \arg\max_{t \in V_v} d_0(t)$
  (farthest tract from $\mathrm{seed}_0$; ties broken by sorted index).
\end{enumerate}
\end{definition}

K-farthest initialisation maximises the initial distance between the two seeds,
preventing the degenerate configuration where both seeds begin in the same
geographic region.
For path-like subgraphs (rural county chains), k-farthest trivially places
seeds at the two endpoints of the path; Voronoi assignment then bisects the
path at the midpoint, and convergence occurs in 1--2 iterations.
For complex urban subgraphs with multiple geographic centres, the farthest-pair
seeds maximise the chance that each seed captures a distinct region.

\begin{remark}
The sorted-index tie-breaking in both the first seed selection
($s_v \bmod m$, which maps the hash to a position in the sorted array) and
the $\arg\max$ ensures that the same $\mathtt{base\_seed}$ and
$\mathtt{node\_path}$ always produce the same seed pair, regardless of the
order in which tracts are processed by the surrounding parallelism.
This reproducibility invariant is tested at L0.
\end{remark}

\subsection{Voronoi Assignment and Medoid Update}

\begin{definition}[Graph Voronoi Assignment]
\label{def:voronoi-assign}
Given seeds $(\mathrm{seed}_0, \mathrm{seed}_1)$, the Voronoi assignment
$a: V_v \to \{0,1\}$ assigns each tract to its nearest seed by BFS distance:
\[
  a(t) = \arg\min_{j \in \{0,1\}} d_{\mathrm{BFS}}(t, \mathrm{seed}_j),
\]
with ties broken in favour of district 0.
\end{definition}

Multi-source BFS computes both $d_{\mathrm{BFS}}(\cdot, \mathrm{seed}_0)$ and
$d_{\mathrm{BFS}}(\cdot, \mathrm{seed}_1)$ simultaneously in a single BFS pass
seeded from both seeds, costing $O(m)$ per assignment step.

\begin{definition}[Approximate Graph-Distance Medoid]
\label{def:medoid}
Let $D_j = \{t \in V_v : a(t) = j\}$ be the tracts assigned to district $j$.
Let $\mathrm{probe}_j \subseteq D_j$ be a uniform random sample of at most
$n_{\mathrm{probe}} = 50$ tracts from $D_j$ (or all tracts if $|D_j| \le 50$).
The \emph{approximate medoid} of $D_j$ is
\[
  \mathrm{medoid}_j = \arg\min_{i \in \mathrm{probe}_j}
  \sum_{t \in D_j} d_{\mathrm{BFS}}(i, t).
\]
\end{definition}

The probe set limits the number of BFS evaluations per district per iteration
to at most 50, making each medoid update $O(50 \times m)$ in the worst case.
For $|D_j| \le 50$, the probe set equals $D_j$ and the medoid is exact.

\subsection{Formal Pseudocode}

Algorithm~\ref{alg:cvd} gives the complete \cvd{} bisection procedure.

\begin{algorithm}[t]
\caption{\cvd{}-\textsc{Bisect}($G_v$, $k_L$, $k_R$, $\mathtt{base\_seed}$,
  $\mathtt{node\_path}$, $\mathtt{cvd\_iters}$, $\mathtt{balance\_tolerance}$)}
\label{alg:cvd}
\begin{algorithmic}[1]
\State $m \gets |V_v|$
\If{$m < 2$}
  \textbf{return} degenerate assignment \Comment{fallback: $\le 1$ tract}
\EndIf
\State $s_v \gets \mathrm{CVD\text{-}Node\text{-}Seed}(\mathtt{base\_seed}, \mathtt{node\_path})$
       \Comment{Definition~\ref{def:cvd-seed}}
\State $\mathrm{sorted} \gets \mathrm{sort}(V_v)$
\Comment{deterministic ordering}
\State $\mathrm{seed}_0 \gets \mathrm{sorted}[s_v \bmod m]$
\State $d_0 \gets \mathrm{BFS\text{-}Distances}(\mathrm{seed}_0, G_v)$
\State $\mathrm{seed}_1 \gets \arg\max_{t \in V_v} d_0[t]$
       \Comment{k-farthest: Definition~\ref{def:kfarthest}}
\State $\mathrm{seeds} \gets (\mathrm{seed}_0, \mathrm{seed}_1)$
\State $\mathrm{prev\_seeds} \gets (-1, -1)$
       \Comment{sentinel: no previous seeds}
\For{$\mathrm{iter} = 0$ \textbf{to} $\mathtt{cvd\_iters} - 1$}
  \State $a \gets \mathrm{VoronoiAssign}(\mathrm{seeds}, G_v)$
         \Comment{multi-source BFS: Definition~\ref{def:voronoi-assign}}
  \For{$j \in \{0, 1\}$}
    \State $D_j \gets \{t : a(t) = j\}$
    \State $\mathrm{probe}_j \gets \mathrm{Sample}(D_j, \min(50, |D_j|))$
    \State $\mathrm{seeds}[j] \gets \arg\min_{i \in \mathrm{probe}_j}
           \sum_{t \in D_j} d_{\mathrm{BFS}}(i, t)$
           \Comment{approx.\ medoid: Definition~\ref{def:medoid}}
  \EndFor
  \If{$\mathrm{seeds} = \mathrm{prev\_seeds}$}
    \textbf{break}
    \Comment{convergence: seeds unchanged}
  \EndIf
  \State $\mathrm{prev\_seeds} \gets \mathrm{seeds}$
\EndFor
\State $a \gets \mathrm{VoronoiAssign}(\mathrm{seeds}, G_v)$
       \Comment{final assignment from converged seeds}
\State $(S, \bar{S}) \gets \mathrm{Rebalance}(a, G_v, \mathbf{p}, \mathtt{balance\_tolerance})$
       \Comment{200-iter boundary-swap}
\State \textbf{return} $(S, \bar{S})$
\end{algorithmic}
\end{algorithm}

\subsection{Three Key Design Decisions}

\paragraph{Decision 1: K-farthest seed placement.}
Na\"{i}ve seed initialisation (e.g., placing both seeds at the first two tracts
in sorted order) would typically start both seeds in the same geographic corner
of the subgraph, requiring many iterations to separate.
K-farthest initialisation places seeds at the two most geographically distant
tracts in the subgraph, maximising coverage from the first iteration.
This reduces the typical iteration count from $\sim 20$ (random initialisation)
to $\sim 10$ (k-farthest) on NC-scale subgraphs.

\paragraph{Decision 2: Medoid, not centroid.}
In continuous Euclidean space, the centroid of a cluster is the well-defined
population-weighted mean position.
In discrete graph space, the mean of a set of tract indices is a fraction
--- a position that may not correspond to any real tract, and that cannot
be used as a BFS source in subsequent Voronoi assignments.
The graph-distance medoid avoids this problem by requiring the new seed to be
an actual tract in its district.
Medoids are also more robust to outliers: the medoid minimises absolute
distance (L1 sense in graph space), whereas the centroid minimises squared
distance (L2 sense).

\paragraph{Decision 3: Post-hoc rebalance.}
Voronoi assignment by BFS distance does not guarantee population balance.
In a state with highly unequal tract populations, the geometrically central
tract may be much more or less populous than others at equal hop count.
The 200-iteration boundary-swap rebalance (consistent with the existing
\texttt{run\_multiscale} rebalance) corrects population imbalance by moving
boundary tracts from the overweight district to the underweight district,
one tract at a time, while preserving contiguity.
Without this step, \cvd{} plans would frequently violate the federal
$\pm 0.5\%$ population balance tolerance.

\subsection{Convergence Analysis}

\begin{proposition}[Monotone Energy Descent]
\label{prop:convergence}
Let $\mathcal{E}(\mathrm{seeds}) = \sum_{j} \sum_{t \in D_j} d_{\mathrm{BFS}}(t, \mathrm{seeds}[j])$
be the total graph-distance quantisation energy.
Each complete iteration (Voronoi assignment + exact medoid update) satisfies
$\mathcal{E}(\mathrm{seeds}^{(i+1)}) \le \mathcal{E}(\mathrm{seeds}^{(i)})$.
\end{proposition}

\begin{proof}
Two observations:
(a) Voronoi assignment minimises $\mathcal{E}$ over all assignments given
    fixed seeds: for each tract $t$, assigning to the nearer seed (by
    $d_{\mathrm{BFS}}$) reduces or maintains $\sum_t d_{\mathrm{BFS}}(t, \mathrm{seeds}[a(t)])$.
(b) Medoid update minimises $\mathcal{E}$ over all seed choices given fixed
    assignment: by definition, $\mathrm{medoid}(D_j)$ minimises
    $\sum_{t \in D_j} d_{\mathrm{BFS}}(i, t)$ over $i \in D_j$.
Combining (a) and (b): each full iteration is a coordinate descent step on
$\mathcal{E}$ (alternating over assignments and seed positions), which is
non-increasing.
Since $\mathcal{E} \ge 0$ and takes values in a finite set (there are finitely
many tract pairs), the sequence $\mathcal{E}(\mathrm{seeds}^{(i)})$ converges
in finitely many steps.
\end{proof}

\begin{remark}
Proposition~\ref{prop:convergence} applies to the exact medoid.
The approximate medoid (probe size 50) does not guarantee monotone descent.
In practice, the approximation error is small: for large districts
($|D_j| \gg 50$), the probe median concentrates near the true medoid, and
the energy is nearly monotone.
For small districts ($|D_j| \le 50$), the probe is exact and the guarantee holds.
\end{remark}

\begin{remark}[Approximate medoid error bound]
An approximate medoid with probe size $p$ has expected relative error bounded by
$O(1/\sqrt{p})$ for the objective function value (the sum of BFS distances).
For $p = 50$: the probability that the true medoid lies outside the sampled set
is $\le (m-p)/m \le 1 - 50/m$.
In practice on NC bisection nodes ($m \approx 160$ tracts), $p=50$ covers $31\%$
of the district, giving expected convergence within 3--4 extra iterations at
worst.
We empirically observe that convergence is unaffected by probe size for
$p \ge 30$ on all test states.
\end{remark}

\begin{remark}[Path-like graphs.]
For subgraphs that are paths (a sequence of tracts $t_0 - t_1 - \cdots - t_{m-1}$),
k-farthest initialisation places $\mathrm{seed}_0$ and $\mathrm{seed}_1$
at the two endpoints.
The Voronoi assignment bisects the path at the midpoint.
The medoid of each half is the midpoint of that half.
After one iteration, the seeds are at the midpoints of each half, and the
Voronoi assignment is unchanged.
Convergence occurs at iteration 2.
\end{remark}

\subsection{Computational Complexity}

\begin{proposition}[Runtime Complexity]
\label{prop:runtime}
\cvd{} bisection at a node with $m$ tracts requires:
\begin{itemize}
\item One BFS for k-farthest initialisation: $O(m)$.
\item Per iteration: two multi-source BFS passes ($O(m)$ each) for
  Voronoi assignment, plus at most $2 \times 50$ single-source BFS
  evaluations ($O(m)$ each) for medoid update.
\item Total per iteration: $O(52 \times m) = O(m)$.
\item Over $n_{\mathrm{iter}}$ iterations: $O(n_{\mathrm{iter}} \times m)$.
\item Post-hoc rebalance: $O(200 \times m)$ for boundary-swap iterations.
\end{itemize}
Total: $O((n_{\mathrm{iter}} + 200) \times m)$.
For the default $\mathtt{cvd\_iters} = 20$, this is $O(220 \times m) = O(m)$
per bisection node.
\end{proposition}

Comparing to \metis{} ($O(m \log m)$ per node) and \sa{} ($O(\mathtt{steps} \times m^{3/2})$
per node in the planar-graph regime): \cvd{} is $O(m)$ per node,
asymptotically faster than both on large subgraphs.
In practice, \cvd{} and \metis{} have similar wall-clock times because the
BFS constant in \cvd{} is larger than the multilevel coarsening constant
in \metis{}.

\subsection{CLI Invocation}

\cvd{} bisection is exposed via \texttt{--structure centroidal-voronoi}
in the \texttt{redist} binary:

\begin{verbatim}
redist state --state NC --year 2020 \
  --structure centroidal-voronoi --cvd-iters 20
\end{verbatim}

With explicit metric and seeding options (Phase~2 preview):
\begin{verbatim}
redist state --state NC --year 2020 \
  --structure centroidal-voronoi --cvd-iters 20 \
  --cvd-metric graph-distance --cvd-seed-init kmeans++
\end{verbatim}

The YAML configuration format:
\begin{verbatim}
algorithm:
  structure: centroidal-voronoi
  weights: geographic
  search: convergence
  cvd_iters: 20
  cvd_metric: graph-distance
  cvd_seed_init: kmeans++
  convergence_threshold: 600
  balance_tolerance: 0.5
workers: 8
\end{verbatim}

\subsection{Audit Chain}

Every \cvd{} run appends structured metadata to
\texttt{runs/\{label\}/\{year\}/index.json}:

\begin{verbatim}
"structure": "centroidal-voronoi",
"n_iter": 20,
"metric": "graph-distance",
"seed_init": "kmeans++",
"base_seed": 12345678,
"cvd_seed": 987654321,
"iterations_to_convergence": 14,
"convergence_warning": false,
"final_ec": 2847,
"final_pp_mean": 0.241
\end{verbatim}

\texttt{iterations\_to\_convergence} records how many iterations were needed
before seeds stabilised.
If \texttt{iterations\_to\_convergence == n\_iter}, the algorithm did not
converge within the budget; \texttt{convergence\_warning} is set to
\texttt{true} as a non-fatal diagnostic.
An auditor who knows \texttt{base\_seed} can recompute
$\mathtt{cvd\_seed} = \mathrm{SHA\text{-}256}(\texttt{"CVD\_INIT\_"} \| \mathtt{base\_seed})$
and verify full reproducibility of the plan.
